% Copyright 2026 The terCAD team. All rights reserved.
% Use of this content is governed by a CC BY-NC-ND 4.0 license that can be found in the LICENSE file.

\subsection{NativePHP}
\markboth{Bootcamping}{NativePHP} \label{nativephp}

NativePHP is an open-source project that converts Laravel applications into native ones. Created
by Marcel Pociot and Simon Hamp, it was presented at Laracon US 2023 for desktop applications
(on top of Electron). Later, the nativephp/mobile package (maintained by Simon Hamp and Shane
Rosenthal) brought the same concept to iOS and Android. Rather than rewriting an application in
Swift or Kotlin, NativePHP packs the PHP interpreter, Laravel code, and dependencies into a
native shell. It also provides access to device capabilities, such as the camera, biometrics,
notifications, and secure storage, through familiar Laravel facades.

\paragraph{Embedded Runtime} PHP is compiled as a shared library (Embed SAPI, thread-safe mode)
and loaded by a native host. On Android, it is Kotlin with Jetpack Compose, and on iOS, it is
Swift with SwiftUI. During startup, the Laravel bundle (shipped as an archive within the
application) is extracted into the application storage. Meanwhile, sessions, the cache, logs,
and the SQLite database are kept in a separate persistent directory that survives application
updates. No network server is involved; requests are passed directly to PHP in the same process.

\paragraph{Rendering Modes} NativePHP offers two ways to draw an interface:

\begin{itemize}
  \item WebView mode (\emph{request / response}) -- every tap or form submission boots PHP,
  renders an HTML page, and returns it to an embedded browser. It's stateless, and an existing
  Laravel application (Blade, Livewire) works almost unchanged.

  \item Native UI mode (\emph{long-lived event loop}) -- PHP starts once and keeps a component
  alive, while its Blade view describes a tree of native elements (real Material and SwiftUI
  buttons, lists, inputs). User actions are sent back to PHP as events, the component updates
  its state, and the screen is re-rendered.
\end{itemize}

\paragraph{Bridge} connects PHP with the native layer in both directions. PHP calls native
functions, such as showing a toast, reading a secure value, or taking a photo, and the native
side reports the results as Laravel events, such as \q{PhotoTaken}, \q{ButtonPressed}, or
\q{Completed}. Plugins deliver platform-specific features and are regular Composer packages with
Kotlin and Swift sources.


\subsubsection{Installing Framework}

NativePHP requires PHP 8.3+ (our project uses 8.4), Composer, Node.js, and native toolchains:
Android Studio (SDK and emulator) for Android, and Xcode (macOS only) for iOS. It's installed
into an existing Laravel application:

\begin{lstlisting}[language=terminal]
$ composer require nativephp/mobile
$ php artisan native:install
$ npm run build # compile artifacts
$ php artisan native:run android
\end{lstlisting}

\noindent The application identity is defined in \q{.env}-file and mapped by
\q{config/nativephp.php}:

\begin{lstlisting}[language=terminal]
NATIVEPHP_APP_ID=com.tercad.fantasyland
NATIVEPHP_APP_VERSION=1.0.0
NATIVEPHP_APP_VERSION_CODE=1
NATIVEPHP_START_URL=/
NATIVEPHP_DEVELOPMENT_TEAM=XXXXXXXXXX # Apple Team ID
\end{lstlisting}

\noindent Since the whole project is shipped to a device, the configuration also declares which
\q{.env}-keys and files should be removed from the bundle (\q{cleanup\_env\_keys},
\q{cleanup\_exclude\_files}), and Android build options (minimal SDK, minification, obfuscation).


\subsubsection{Running and Debugging}

The \q{native:*}-commands cover the daily development cycle:

\begin{lstlisting}[language=terminal]
# build, install, hot reload
$ php artisan native:run ios --watch
# start an emulator
$ php artisan native:emulator android
# stream logs (Android app)
$ php artisan native:tail
# open Android Studio / Xcode
$ php artisan native:open android
# collect environment details
$ php artisan native:debug
\end{lstlisting}

\noindent NativePHP enables WebView debugging in the native shells
(\q{setWebContentsDebuggingEnabled} on Android, \q{isInspectable} on iOS), so the WebView
content can be inspected by Chrome DevTools (\q{chrome://inspect}) for Android and Safari Web
Inspector for iOS.


\subsubsection{Rendering in WebView}

In the simplest case, our routes and controllers stay the same as for a web application, and a
native screen just embeds them by \q{<native:webview>}-element:

\begin{lstlisting}[language=HTML]
{{-- resources/views/native/home.blade.php --}}
<native:webview php src="{{ $src }}" fullscreen />
\end{lstlisting}

\noindent The \q{php}-attribute indicates that the source is served by the embedded runtime, not
a remote site.


\subsubsection{Building Native Screens}

A native screen is a \q{NativeComponent}, which is a PHP class with public properties
representing state and methods representing actions. A Blade view describes the elements. The
command \q{php artisan native:make Counter} generates both. Components are bound to URIs by
\q{Route::native}, and a layout (\q{NativeLayout}) adds shared navigation, such as a tab bar:

\begin{lstlisting}
// routes/web.php
Route::nativeGroup(NavigationLayout::class, function () {
  Route::native('/shell/home', HomeShell::class); (*@ \stopnumber @*)
});

// app/NativeComponents/HomeShell.php
class HomeShell extends NativeComponent
{
  public string $activeTab = 'home';

  public function mount(): void
  {
    // plugin: srwiez/nativephp-mobile-screen
    Screen::keepAwake();
  }

  public function showChat(): void
  {
    // re-render on state change
    $this->activeTab = 'chat';
  }

  public function render(): View
  {
    return view('native.home',
        ['src' => $this->source()]);
  } (*@ \stopnumber @*)
}

// app/NativeComponents/NavigationLayout.php
class NavigationLayout extends NativeLayout
{
  public function tabBar(NativeComponent $screen): ?TabBar
  {
    return (new TabBar)
      ->dark()
      ->add(Tab::action('Home', icon: 'home')->press('showHome'))
      ->add(Tab::action('Chat', icon: 'chat_bubble')->press('showChat'));
  }
}
\end{lstlisting}

\noindent Elements are rendered by platform primitives: the core package provides layout and
basic controls (\q{column}, \q{row}, \q{text}, \q{button}, \q{toggle}, \q{bottom-sheet}, etc.),
and the \q{nativephp/mobile-ui}-package adds richer ones (\q{select}, \q{pager},
\q{outlined-text-input}, \q{modal}, etc.). Layouts accept utility classes
(\q{class="gap-2 p-4"}), controls are styled by a theme, and events are bound by \q{@press} and
\q{native:model}:

\newpage
\begin{lstlisting}[language=HTML]
<native:column class="gap-2 p-4">
  <native:text :fontSize="24">
    Count: {{ $count }}
  </native:text>
  <native:outlined-text-input label="Name" native:model.live="name" />
  <native:button variant="primary" @press="increment">
    Increment
  </native:button>
</native:column>
\end{lstlisting}


\subsubsection{Managing Component State}

Attributes fine-tune the behavior of components and their methods:

\begin{itemize}
  \item \q{\#[Computed]} -- memoizes a method result per render frame (\q{\$this->revenue});
  \item \q{\#[Locked]} -- prevents a property from being changed by two-way bindings (\q{native:model});
  \item \q{\#[Lazy]} -- shows a placeholder while a slow \q{mount()} is executed;
  \item \q{\#[Poll]} -- re-runs a method periodically;
  \item \q{\#[On]} and \q{\#[OnNative]} -- listen for application and device events.
\end{itemize}

\begin{lstlisting}
#[Lazy]
class Profile extends NativeComponent
{

  #[Locked]
  public int $heroId;

  #[Computed(persist: true)]
  public function stats(): array
  {
    return Hero::find($this->heroId)->stats();
  }
}
\end{lstlisting}


\subsubsection{Accessing Device Capabilities}

Device APIs are exposed as facades in the \q{Native\textbackslash Mobile\textbackslash Facades}
-namespace:

\begin{lstlisting}
use Native\Mobile\Facades\{Camera, Device, Dialog, Haptics, SecureStorage};

Dialog::toast('New message received');
Dialog::alert('Exit', 'Leave the fight?', ['Yes', 'No']);
Dialog::share('Fantasyland', 'Join me!', 'https://fantasyland.ru');
Haptics::vibrate();
Device::getBatteryInfo();

// Keychain (iOS) / Keystore (Android)
SecureStorage::set('token', $token);
$token = SecureStorage::get('token');

// Asynchronous call, completed by an event
Camera::getPhoto();
\end{lstlisting}

\noindent Results of asynchronous calls arrive as events, which components handle by \q{\#[OnNative]}:

\begin{lstlisting}
use Native\Mobile\Attributes\OnNative;
use Native\Mobile\Events\Camera\PhotoTaken;
#[OnNative(PhotoTaken::class)]
public function photoTaken(string $path): void
{
  $this->avatar = $path;
}
\end{lstlisting}

\noindent Helpers \q{isIos()}, \q{isAndroid()}, \q{isMobile()}, and \q{isDark()} allow to adapt
the behavior per platform and theme.


\subsubsection{Running in Background}

Laravel queues and the scheduler also work on the device, so jobs, such as our chat listener,
are dispatched in the usual way. For heavy computations within a native screen, \q{AsyncTask}
executes a static closure on a separate PHP thread and returns the result to the component:

\begin{lstlisting}
public function generateReport(): void
{
  $this->loading = true;

  AsyncTask::dispatch(static fn() => Report::build()->toArray())
    ->finished(function (array $result) {
      $this->report = $result;
      $this->loading = false;
    })
    ->failed(fn(Throwable $e) => $this->error = $e->getMessage());
}
\end{lstlisting}


\subsubsection{Extending by Plugins}

Plugins are Composer packages with PHP facades and native (Kotlin/Swift) implementations. For
security reasons, only the plugins listed in \q{NativeServiceProvider::plugins()} are compiled
into a build. This prevents transitive dependencies from registering themselves silently:

\begin{lstlisting}[language=terminal]
$ composer require srwiez/nativephp-mobile-screen
$ php artisan vendor:publish --tag=nativephp-plugins-provider
$ php artisan native:plugin:register srwiez/nativephp-mobile-screen
$ php artisan native:plugin:list
\end{lstlisting}

\begin{lstlisting}
public function plugins(): array
{
  return [
    \Native\Mobile\UI\NativeUIServiceProvider::class,
    \SRWieZ\NativePHP\Mobile\Screen\ScreenServiceProvider::class,
  ];
}
\end{lstlisting}

\noindent Own plugins are scaffolded by \q{native:plugin:create} and checked by
\q{native:plugin:validate}.


\subsubsection{Packaging for Stores}

Release builds are signed and packaged into an Android App Bundle (\q{.aab}) or an iOS archive (\q{.ipa}):

\begin{lstlisting}[language=terminal]
# bump version: 1.0.0 -> 1.0.1
$ php artisan native:release patch
# generate a keystore
$ php artisan native:credentials android
$ php artisan native:package android --build-type=bundle --keystore=release.jks --key-alias=fantasyland
$ php artisan native:package ios --export-method=app-store
$ php artisan native:check-build-number
\end{lstlisting}

\noindent The result is uploaded to Google Play Console, App Store Connect (TestFlight), or
alternative stores, as we'll discover in the distribution chapters.


\subsubsection{Asserting in Tests}

Native components are tested without a device: \q{Native::test()} mounts a component, simulates
user actions, and checks the rendered tree, while \q{Native::fakeBridge()} replaces responses of
native calls:

\begin{lstlisting}
use Native\Mobile\Testing\Native;

Native::test(HomeShell::class)
  ->call('showChat')
  ->assertSet('activeTab', 'chat')
  ->assertElement('webview');

Native::fakeBridge()->respondTo('Geolocation.GetCurrentPosition', [
  'latitude' => 55.75, 'longitude' => 37.61,
]);

Native::visit('/shell/home')->assertScreen(HomeShell::class);
\end{lstlisting}

\noindent Pest tests for screens are generated by \q{php artisan native:make-test}, and
\q{AsyncTask::fake()} verifies dispatched background tasks (\q{assertDispatched},
\q{assertDispatchedTimes}).
